\scanpage{129}
\mspage{104}
% letterhead
docteur alain lambin-dostromon\\
lo marchadiol\\
turenne 19500 meyssac\\
tél. : (55) 85-90-23\\
sur rendez-vous\\

\margin{\unsure{XXX}}

Octobre 1985\\
à\qquad Mr SOLE\\
Directeur Adjoint CPAM\\

% slanted
\margin{COPIE}

Monsieur le Directeur,\\

Je reçois bien votre courrier du 2 septembre dernier relatif à la ``campagne de vaccination contre la grippe en faveur des personnes âgées''.

Je ne puis me soustraire au devoir que m'impose le code de déontologie, qui me fait obligation de ne pas garder secret ce que je sais.

Vous écrivez que ``cette expérience (de vaccination sur 3 ans) a permis une amélioration de la couverture vaccinale des personnes âgées''. Ceci est une évidence, mais ne signifie pas l'amélioration de l'état de santé de la population en question.

Je me dois de vous informer que tout vaccin allopathique a pour but d'augmenter les défenses d'un organisme déficient. D'où augmentation des gamma globulines ; \unsure{mais}... celles ci restent élevées devant l'agression vaccinale, et l'on crée ainsi un terrain tuberculinique, au sens de l'Homéopathie et de la Médecine Quantique.

Ceci va bientôt s'agraver d'une baisse des sérines, correspondant à une sycose. L'association des deux phénomènes crée une tuberculinosycose, ou cancérinisme, qui fait le lit.

- des cancers (vieillissement cellulaire et défaut de transfert des informations).

- des scléroses d'organes (insuffisance rénale, cardiaque, hépatique, diabète, artérite, artériosclérose, etc...) pour 85 \% des cas et aux

- scléroses nerveuses (débilité mentale, schizophrénie, sclérose en plaque, maladie de charcot, sclérose latérale, \unsure{amyotrophique}, etc...) pour 15 \% des cas.

Il existe d'autres évolutions inverses, sous vaccination, chez les sujets plus jeunes, avec hypo gamma globulinémie et hyper-sérine.

Ces conclusions, issues de nos travaux au sein de la Société Internationale de Médecine Quantique , et de l'utilisation de bilans protéiques informatisés, avec un recul d'au moins 10 ans pour nous même.

Je dois avouer que mon attention a été attirée sur les pathologies iatrogènes dues aux vaccins par le fait que je suis moi même une victime du vaccin antipolio injectable, réputé totalement atoxique... par les laboratoires qui le vendent !!! Et qui sont ceux que vous citez.

% letterhead
docteur alain lambin-dostromon\\
lo marchadiol\\
turenne 19500 meyssac\\
tél. : (55) 85-90-23\\
sur rendez-vous\\

Je ne critique pas leurs recherches, mais leurs méthodes commerciales.

Je me permets de vous citer mon expérience au Cabinet médical :

La protection par les vaccins que vous citez est assez correcte contre la grippe sauf :

- que les suites graves citées plus haut n'ont pas été mises en rapport avec la vaccination, dans l'esprit des allopathes. Peut être parce que ceux ci ne voient plus les maladies graves dont je parle.

- que les grippes graves déclanchées par le vaccin ne sont pas négligeables ni en nombre, ni quand aux conséquences vitales parfois, je suis surpris, aussi, que vous ne connaissiez pas les protections homéopathiques par Influenzinum ou Oscillococcinum selon l'âge du sujet.

Infiniment plus efficace et sans danger du type de ceux cités plus haut. Mon expérience de dix ans le prouve aisément, et se vérifie par tout médecin de bonne volonté.

Produits moins chers, pas dangereux, faciles d'emploi, baissans la consommation médicale. Il me semblerait plus logique de les intégrer à la liste des produits que vous citez.

Ainsi les personnes âgées auraient le choix, de leur thérapeutique, conformément aux droits de l'Homme, dont on parle d'autant plus qu'on les baffoue perfidement.

Pourriez vous avoir l'amabilité de me répondre. Ces courriers pouvant éventuellement être publiés.

Vous comprendrez aussi que, médicalement parlant, je ne puisse adhérer à une campagne vaccinale qui risque de créer encore plus de malades âgées, et ruiner un peu plus mes finances qui les alimentent, et le soignent.

Dans l'attente de votre réponse, je vous prie d'agréer, Monsieur\\
le Directeur, mes salutations les plus sincères.\\

\handwritten{\unsure{A|AL}}\\

PS : à votre disposition pour vous éclairer mieux sur la question.
% archive sheet number: 129
